\begin{figure}[H]
  \centering
  \begin{tikzpicture}
    \begin{axis}[
      width=0.82\textwidth,
      height=0.52\textwidth,
      xlabel={Nồng độ kali (mg/L)},
      ylabel={Cường độ phát xạ},
      xmin=0, xmax=26,
      ymin=0, ymax=27,
      xtick={0,5,10,15,20,25},
      ytick={0,5,10,15,20,25},
      grid=both,
      major grid style={gray!35},
      minor grid style={gray!15},
      legend style={at={(0.97,0.04)},anchor=south east,font=\small},
      legend cell align={left}
    ]
      \addplot[
        only marks,
        mark=*,
        mark size=2.3pt,
        blue!75!black
      ] coordinates {(0,0) (5,5) (10,10) (15,16) (20,21) (25,25)};
      \addlegendentry{Số liệu chuẩn}
      \addplot[red!75!black, very thick, domain=0:25, samples=100] {1.022857142857*x + 0.047619047619};
      \addlegendentry{Hồi quy tuyến tính}
      \node[
        anchor=north west,
        draw=black,
        fill=white,
        fill opacity=0.92,
        text opacity=1,
        align=left,
        font=\small
      ] at (axis cs:1,25.8) {$y = 1{,}0229x + 0{,}0476$\\$R^2 = 0{,}9976$};
    \end{axis}
  \end{tikzpicture}
  \caption{Đường hồi quy tuyến tính của dãy chuẩn kali}
  \label{fig:duong-hoi-quy-kali}
\end{figure}
